% =============================================================
% บทที่ 1 - บทนำ
% โครงงานประเภท B (งานวิจัย) — เน้น Research Gap และ Motivation
% ขอบเขตระบุ Dataset ตัวแปร และขอบเขตการทดลอง
% =============================================================

% -----------------------------------------------------------
\section{ความเป็นมาและความสำคัญของปัญหา}
% -----------------------------------------------------------

ในปัจจุบัน การค้นหาสินค้าบนร้านค้าออนไลน์เป็นกิจกรรมที่ผู้บริโภคทำเป็นประจำ ระบบค้นหาที่ร้านค้าออนไลน์ใช้กันอยู่ส่วนใหญ่ทำงานโดยจับคู่คำที่ผู้ใช้พิมพ์กับคำที่ปรากฏในชื่อสินค้า วิธีนี้ได้ผลดีเมื่อผู้ใช้รู้อยู่แล้วว่าต้องการสินค้าชื่ออะไร แต่ได้ผลไม่ดีเมื่อผู้ใช้บอกความต้องการเป็นภาษาธรรมชาติ เช่น ระบุเงื่อนไขหลายข้อพร้อมกัน บอกงบประมาณ หรือใช้คำเรียกที่ต่างจากคำที่ร้านค้าใช้

แนวทางที่ได้รับความสนใจในระยะหลังคือการนำแบบจำลองภาษาขนาดใหญ่มาช่วยตอบ โดยให้ค้นข้อมูลจากคลังสินค้าก่อนแล้วจึงเขียนคำตอบจากสิ่งที่ค้นได้ เรียกว่าการสร้างคำตอบโดยอาศัยการค้นคืน หรือ Retrieval-Augmented Generation \cite{lewis2020rag} วิธีนี้ทำให้ระบบตอบด้วยข้อมูลสินค้าจริงแทนการแต่งขึ้นเอง

อย่างไรก็ตาม โครงสร้างพื้นฐานของวิธีนี้ทำงานเป็นท่อเส้นตรง คือรับคำถาม ค้นหนึ่งครั้ง แล้วเขียนคำตอบ จุดอ่อนที่ตามมาคือเมื่อการค้นครั้งเดียวนั้นหาของไม่เจอ ระบบไม่มีโอกาสแก้ตัวและยังคงตอบด้วยสินค้าที่ไม่ตรงกับที่ผู้ใช้ขอ ปัญหานี้รุนแรงขึ้นในบริบทของผู้ใช้ชาวไทย เพราะข้อมูลสินค้าในคลังส่วนใหญ่เป็นภาษาอังกฤษ ขณะที่ผู้ใช้พิมพ์คำค้นเป็นภาษาไทย ความคลาดเคลื่อนที่เกิดจากการข้ามภาษาจึงซ้อนทับกับความคลาดเคลื่อนของการค้นอีกชั้นหนึ่ง

เพื่อแก้จุดอ่อนดังกล่าว งานวิจัยจำนวนหนึ่งเสนอให้ระบบตัดสินใจได้เองว่าจะค้นซ้ำหรือไม่ และจะปรับคำค้นอย่างไร \cite{jeong2024adaptive, yan2024crag} แนวทางกลุ่มนี้เรียกรวมว่า Agentic RAG ซึ่งนิยามด้วยการมีชั้นที่แบบจำลองภาษาเป็นผู้ตัดสินใจเองขณะทำงาน ไม่ใช่เพียงทำตามลำดับที่เขียนไว้ล่วงหน้า \cite{singh2025agentic}

\textbf{ช่องว่างที่พบ} คือแม้แนวคิดนี้จะได้รับการยอมรับในทางทฤษฎี แต่หลักฐานเชิงประจักษ์ยังมีข้อจำกัดสามประการ ประการแรก งานที่เสนอสถาปัตยกรรมใหม่มักเปรียบเทียบกับระบบพื้นฐานโดยเปลี่ยนหลายองค์ประกอบพร้อมกัน จึงสรุปไม่ได้ว่าผลที่ดีขึ้นมาจากความสามารถในการตัดสินใจจริงหรือมาจากองค์ประกอบอื่น ประการที่สอง งานเหล่านั้นประเมินบนคำค้นและเอกสารที่เป็นภาษาเดียวกัน ยังไม่มีหลักฐานในสถานการณ์ที่คำค้นกับข้อมูลเป็นคนละภาษา และประการที่สาม งานเหล่านั้นส่วนใหญ่ใช้แบบจำลองภาษาขนาดใหญ่ผ่านบริการภายนอก จึงยังไม่ทราบว่าแบบจำลองขนาดที่ทำงานได้บนเครื่องคอมพิวเตอร์ส่วนบุคคลจะมีความสามารถในการตัดสินใจเพียงพอหรือไม่

ด้วยเหตุนี้ โครงงานนี้จึงมุ่งเปรียบเทียบสถาปัตยกรรมแบบ RAG มาตรฐานกับแบบ Agentic RAG บนระบบแนะนำสินค้าเชิงสนทนาที่รองรับทั้งภาษาไทยและภาษาอังกฤษ โดยออกแบบให้สถาปัตยกรรมที่นำมาเปรียบเทียบใช้ส่วนประกอบร่วมกันทั้งหมดและต่างกันเฉพาะชั้นที่ตัดสินใจ เพื่อให้ตอบได้ว่าความสามารถในการตัดสินใจเองนั้นให้ประโยชน์จริงหรือไม่ และให้ประโยชน์ในสถานการณ์ใด

% -----------------------------------------------------------
\section{วัตถุประสงค์}
% -----------------------------------------------------------

\begin{enumerate}[label=\arabic*)]
    \item เพื่อพัฒนาระบบแนะนำสินค้าเชิงสนทนาที่รับคำถามได้ทั้งภาษาไทยและภาษาอังกฤษ บนคลังสินค้าจริงที่เป็นภาษาอังกฤษ โดยประมวลผลบนเครื่องคอมพิวเตอร์ส่วนบุคคลทั้งหมด
    \item เพื่อเปรียบเทียบประสิทธิภาพระหว่างสถาปัตยกรรมแบบ RAG มาตรฐานกับแบบ Agentic RAG ด้วยการออกแบบการทดลองที่เปลี่ยนตัวแปรทีละตัว
    \item เพื่อศึกษาว่าความสามารถในการตัดสินใจเองของแบบจำลองภาษาขนาดที่ทำงานได้บนเครื่องส่วนบุคคล ส่งผลต่อคุณภาพการแนะนำสินค้าในสถานการณ์ใดบ้าง
    \item เพื่อสร้างและเผยแพร่ชุดทดสอบพร้อมเฉลยสำหรับการค้นหาสินค้าสองภาษา ให้ผู้อื่นนำไปวัดผลซ้ำได้
\end{enumerate}

% -----------------------------------------------------------
\section{ขอบเขตและข้อจำกัดของโครงงาน}
% -----------------------------------------------------------

\subsection{ขอบเขตด้านข้อมูล}

\begin{enumerate}[label=\arabic*)]
    \item ใช้ข้อมูลสินค้าจริงจากคลัง Amazon Reviews 2023 \cite{hou2024amazon} หลังผ่านการล้างข้อมูลเหลือ 81{,}264 รายการ ครอบคลุมสินค้าหลายหมวด
    \item ข้อมูลสินค้าทั้งหมดเป็นภาษาอังกฤษ ไม่มีรายการใดที่ชื่อสินค้าเป็นภาษาไทย
    \item ไม่มีการเพิ่มหรือปรับปรุงข้อมูลสินค้าระหว่างการทดลอง คลังข้อมูลถูกตรึงไว้ตลอดการวัดผล
\end{enumerate}

\subsection{ขอบเขตด้านระบบที่พัฒนา}

\begin{enumerate}[label=\arabic*)]
    \item พัฒนาสถาปัตยกรรม 3 แบบเพื่อนำมาเปรียบเทียบ ได้แก่ แบบ RAG มาตรฐานที่ไม่มีชั้นตัดสินใจ แบบที่แบบจำลองภาษาตัดสินใจเฉพาะการเลือกการกระทำ และแบบที่แบบจำลองภาษาตัดสินใจทั้งการเลือกการกระทำและการค้นซ้ำ
    \item ทั้งสามแบบใช้ส่วนประกอบร่วมกันทั้งหมด ได้แก่ ส่วนความจำ ส่วนแปลภาษา ส่วนค้นคืน ส่วนจัดอันดับ และส่วนเขียนคำตอบ ต่างกันเฉพาะชั้นที่ตัดสินใจเท่านั้น
    \item ระบบรับคำถามเป็นข้อความ สนทนาได้หลายรอบ และจดจำความต้องการของผู้ใช้ข้ามรอบสนทนา
    \item ระบบแนะนำสินค้าครั้งละไม่เกิน 3 รายการ และต้องตอบว่าไม่พบเมื่อคลังไม่มีสินค้าที่ตรงกับที่ขอ
    \item ไม่ได้พัฒนาส่วนติดต่อผู้ใช้แบบเว็บหรือแอปพลิเคชันบนมือถือ เนื่องจากโครงงานนี้มุ่งตอบคำถามเชิงเปรียบเทียบสถาปัตยกรรม ไม่ใช่การส่งมอบผลิตภัณฑ์
\end{enumerate}

\subsection{ขอบเขตด้านการทดลองและการวัดผล}

\begin{enumerate}[label=\arabic*)]
    \item สร้างชุดทดสอบจำนวน 140 คำค้น โดยยืมข้อความคำค้นจริงของผู้ใช้จากชุดข้อมูลมาตรฐานที่ตีพิมพ์แล้ว \cite{reddy2022esci, chen2022wands} และแปลเป็นภาษาไทยคู่กัน
    \item แบ่งชุดทดสอบเป็น 3 กลุ่มตามสิ่งที่คลังมี ได้แก่ กลุ่มที่คลังมีสินค้าตรง กลุ่มที่คลังไม่มีสินค้าตรง และกลุ่มที่มีเงื่อนไขงบประมาณ
    \item ติดเฉลยความตรงของสินค้าด้วยมือตามเกณฑ์มาตรฐาน 4 ระดับ \cite{reddy2022esci} และติดเฉลยความซื่อตรงของคำตอบตามเกณฑ์ 3 ระดับ
    \item วัดผลด้วย nDCG@3 \cite{jarvelin2002ndcg} ร่วมกับตัวชี้วัดการตัดสินใจแนะนำหรือปฏิเสธ การทำตามงบประมาณ ความซื่อตรงของคำตอบ และต้นทุนด้านเวลา
    \item ทดสอบนัยสำคัญทางสถิติด้วยการทดสอบจับคู่แบบสลับเครื่องหมายที่นับครบทุกกรณี \cite{smucker2007comparison}
    \item ประมวลผลบนเครื่องคอมพิวเตอร์ส่วนบุคคลเครื่องเดียวตลอดการทดลอง ไม่เรียกใช้บริการแบบจำลองภาษาผ่านเครือข่าย
\end{enumerate}

\subsection{ข้อจำกัดของโครงงาน}

\begin{enumerate}[label=\arabic*)]
    \item แบบจำลองภาษาที่ใช้มีขนาด 8 พันล้านพารามิเตอร์ ซึ่งเป็นขนาดใหญ่ที่สุดที่ทำงานได้บนหน่วยความจำของหน่วยประมวลผลกราฟิกขนาด 8 กิกะไบต์ ผลที่ได้จึงสรุปไปถึงแบบจำลองขนาดใหญ่กว่านี้ไม่ได้
    \item คำค้นภาษาไทยในชุดทดสอบได้จากการแปลคำค้นภาษาอังกฤษ ไม่ใช่คำค้นที่ผู้ใช้ชาวไทยพิมพ์ขึ้นเอง ความยากที่เกิดจากการสะกดผิดจึงหายไปบางส่วนเมื่อผ่านการแปล
    \item เฉลยความตรงติดโดยผู้วิจัยซึ่งเป็นผู้พัฒนาระบบที่ถูกประเมินเอง ยังไม่มีผู้ติดป้ายอิสระคนที่สองมาติดซ้ำ
    \item กองที่ติดเฉลยมาจากสินค้าที่ระบบของโครงงานนี้หยิบขึ้นมาเท่านั้น ระบบอื่นที่หยิบสินค้านอกกองจะถูกประเมินต่ำกว่าความเป็นจริง \cite{zobel1998reliable}
    \item ตัวหารของ nDCG ใช้ระบบในอุดมคติแทนกองที่ติดเฉลย ค่าที่ได้จึงนำไปเทียบกับงานวิจัยอื่นที่ใช้ตัวหารแบบอื่นโดยตรงไม่ได้
    \item ชุดทดสอบไม่ครอบคลุมสถานการณ์ที่คลังมีสินค้าใกล้เคียงแต่ไม่ตรง ซึ่งเป็นสถานการณ์ที่การค้นซ้ำน่าจะได้เปรียบมากที่สุด การประเมินจึงอาจต่ำกว่าประโยชน์จริงของการค้นซ้ำ
\end{enumerate}

% -----------------------------------------------------------
\section{ประโยชน์ที่คาดว่าจะได้รับ}
% -----------------------------------------------------------

\begin{enumerate}[label=\arabic*)]
    \item ได้ระบบแนะนำสินค้าเชิงสนทนาที่รับคำถามได้สองภาษาบนคลังสินค้าภาษาอังกฤษ และทำงานได้บนเครื่องคอมพิวเตอร์ส่วนบุคคลโดยไม่ต้องพึ่งบริการภายนอก
    \item ได้หลักฐานเชิงประจักษ์ว่าความสามารถในการตัดสินใจเองของแบบจำลองภาษาให้ประโยชน์ในสถานการณ์ใด และไม่ให้ประโยชน์ในสถานการณ์ใด ซึ่งนำไปใช้ตัดสินใจเลือกสถาปัตยกรรมในงานอื่นได้
    \item ได้ชุดทดสอบการค้นหาสินค้าสองภาษาพร้อมเฉลยที่เผยแพร่สู่สาธารณะ ซึ่งผู้อื่นนำไปวัดผลระบบของตนเองและเทียบกับผลของโครงงานนี้ได้
    \item ได้องค์ความรู้เกี่ยวกับข้อจำกัดของแบบจำลองภาษาขนาดที่ทำงานบนเครื่องส่วนบุคคลได้ เมื่อนำมาใช้ในบทบาทผู้ตัดสินใจภายในระบบค้นคืน
\end{enumerate}
